\documentclass[11pt]{article}
\usepackage[T1]{fontenc}
\usepackage[utf8]{inputenc}
\usepackage{lmodern}
\usepackage{geometry}
\usepackage{amsmath}
\usepackage{amssymb}
\usepackage{amsthm}
\usepackage{booktabs}
\usepackage{hyperref}
\usepackage{listings}
\geometry{margin=1in}

\lstdefinestyle{lean}{
  basicstyle=\ttfamily\small,
  breaklines=true,
  columns=fullflexible,
  keepspaces=true,
  showstringspaces=false,
  frame=single,
  framesep=4pt,
  xleftmargin=2pt,
  aboveskip=8pt,
  belowskip=8pt,
  literate=
    {ℝ}{{$\mathbb{R}$}}1
    {ℕ}{{$\mathbb{N}$}}1
    {ℤ}{{$\mathbb{Z}$}}1
    {λ}{{$\lambda$}}1
    {ε}{{$\varepsilon$}}1
    {Δ}{{$\Delta$}}1
    {κ}{{$\kappa$}}1
    {≤}{{$\leq$}}1
    {≥}{{$\geq$}}1
    {≠}{{$\neq$}}1
    {∧}{{$\wedge$}}1
    {∨}{{$\vee$}}1
    {¬}{{$\neg$}}1
    {→}{{$\rightarrow$}}1
    {↦}{{$\mapsto$}}1
    {∃}{{$\exists$}}1
    {∀}{{$\forall$}}1
    {∈}{{$\in$}}1
    {·}{{$\cdot$}}1
    {²}{{$^2$}}1
    {₀}{{$_0$}}1
    {₁}{{$_1$}}1
    {₂}{{$_2$}}1
    {₃}{{$_3$}}1
    {⟨}{{$\langle$}}1
    {⟩}{{$\rangle$}}1,
}
\lstset{style=lean}

\newtheorem{theorem}{Theorem}
\newtheorem{proposition}{Proposition}
\newtheorem{lemma}{Lemma}
\newtheorem{remark}{Remark}
\theoremstyle{definition}
\newtheorem{definition}{Definition}

\newcommand{\SUN}{\mathrm{SU}(N)}
\newcommand{\Lirr}{\lambda_{\mathrm{irr}}}
\newcommand{\dmax}{d_{\max}}
\newcommand{\Tr}{\mathrm{Tr}}
\newcommand{\Pleq}{P_{\leq \dmax}}

\title{A Machine-Checked Audit of the Extended-Extraction\\
RG Contraction in the Wilson-Lattice / OS Route\\
to a 4D Yang--Mills Mass Gap}
\author{The MeTTapedia formalization effort}
\date{June 2026}

\begin{document}
\maketitle
\sloppy

\begin{abstract}
Ben Goertzel's manuscript \emph{Yang--Mills Mass Gap in 4D: RG Bootstrap,
Polymer Expansion, OS Reconstruction} (v14) proposes to construct a continuum
$\SUN$ Yang--Mills theory as a limit of Wilson lattice gauge theory, organized
around a multiscale renormalization-group (RG) bootstrap, a Koteck\'y--Preiss
(KP) polymer expansion yielding exponential clustering at a fixed physical
scale, and an Osterwalder--Schrader (OS) reconstruction producing a Hamiltonian
with a strictly positive spectral gap. The technical keystone is an
\emph{extended-extraction} RG contraction $\Lambda = C_1\,\Lirr < 1$, with an
advertised recombination constant $C_1 \leq 2224$ and an irrelevant-eigenvalue
factor $\Lirr = 2^{-13}$ (block factor $b=2$, extraction depth $\dmax = 16$), so
that $2224 \cdot 2^{-13} \approx 0.27 < 1$. We report what the MeTTapedia
formalization effort has machine-checked in Lean about this route. We model the
contraction predicate \texttt{HasExtendedExtractionContraction}, verify the
knife-edge arithmetic ($2224 \cdot 2^{-13} < 1$ holds at $\dmax=16$; the
analogous bound \emph{fails} at the nearest even depth $\dmax=14$, where
$2^{11}=2048<2224$), establish genuine finite $\mathbb{Z}_2$ and $\mathbb{Z}_3$
single-link strong-coupling spectral gaps through a finite OS reconstruction, and
\emph{gate the continuum endpoint as open}: the continuum mass gap carries an
explicit open status with no claimed support. We are precise about scope. The
checked arithmetic takes the recombination constant $C_1 = 2224$ as a
\emph{given input}. The load-bearing quantity --- whether the true
extraction-projection operator norm, hence $C_1$, is as small as advertised ---
lies outside the present formalization; we record neutrally that the manuscript's
own written expression for that norm appears to grow rapidly with the extraction
depth $\dmax$, so the small constant is not self-evidently justified, and we flag
this as the key unformalized risk rather than asserting any resolution. The
finite gaps that are checked are real but Clay-irrelevant (abelian, single-link,
strong-coupling, finite). The Yang--Mills mass gap itself remains open; we make
no claim about it.
\end{abstract}

\section{The problem and Ben Goertzel's route}

\subsection{The Millennium problem and the lattice strategy}

Fix the gauge group $G = \SUN$ with $N \geq 2$ and let $\mathbb{R}^4$ be
Euclidean space. The Yang--Mills mass-gap problem asks for the construction of a
nontrivial continuum quantum Yang--Mills theory on $\mathbb{R}^4$ with gauge
group $G$, satisfying the Wightman / Osterwalder--Schrader axioms, whose
reconstructed Hamiltonian has a unique vacuum and a strictly positive spectral
gap on the orthocomplement of the vacuum.

Goertzel's v14 manuscript follows Wilson's lattice regularization. One
discretizes $\mathbb{R}^4$ as a hypercubic lattice $a\mathbb{Z}^4$ with spacing
$a>0$, with link variables $U_{x,\mu}\in G$ on oriented nearest-neighbor links.
The plaquette variable is the parallel transporter around an elementary square,
\[
  U_{x,\mu\nu} \;=\; U_{x,\mu}\,U_{x+a e_\mu,\nu}\,U_{x+a e_\nu,\mu}^{-1}\,U_{x,\nu}^{-1},
  \qquad 1 \leq \mu < \nu \leq 4,
\]
and the Wilson action penalizes deviation of plaquettes from the identity,
\[
  S_{\Lambda,a}(U) \;=\; \beta \sum_{p\subset\Lambda} V(U_p),
  \qquad \beta = \frac{2N}{g^2},
  \qquad V(U) := 1 - \tfrac{1}{N}\mathrm{Re}\,\Tr(U) \in [0,2].
\]
Weak coupling is $\beta \gg 1$. The strategy is: (Step 1) start from the
well-defined finite-dimensional lattice integral; (Step 2) run a multiscale RG
bootstrap controlling the effective theory across all scales; (Step 3)
establish exponential clustering at a fixed physical scale $\ell_*$ via a KP
polymer expansion; (Step 4) transfer clustering from the coarse scale back to
the fine scale with a glue / back-propagation lemma; (Step 5) reconstruct the
continuum Hilbert space, vacuum, and Hamiltonian by OS, with the clustering rate
becoming the spectral gap; (Step 6) verify nontriviality through a nonzero
third cumulant of the action density.

\subsection{The continuum-limit obstruction and extended extraction}

In any RG scheme one must show that the ``irrelevant'' sector (operators of
canonical dimension $>4$) contracts at each step:
\[
  \| K^{\mathrm{irr}}_{j+1}\| \;\leq\; \Lambda\, \|K^{\mathrm{irr}}_{j}\| + (\text{source terms}),
  \qquad \Lambda < 1.
\]
The manuscript identifies a sharp obstruction in the \emph{naive} scheme, where
one extracts only the relevant and marginal operators of dimension $\leq 4$.
Under rescaling, a dimension-$d$ operator scales by $b^{4-d}$, so the leading
irrelevant operators (dimension $6$) contract by $b^{-2} = 0.25$ for $b=2$. But
the RG step also recombines polymers, contributing a combinatorial constant
$C_1$. The manuscript's rigorous worst-case bound is $C_1 \leq 2224$ for $b=2$,
giving
\[
  \Lambda_{\mathrm{naive}} \;=\; C_1\, b^{-2} \;\geq\; 2224 \times 0.25 \;=\; 556 \;>\; 1,
\]
an \emph{expanding} map: initial smallness would grow like $556^{J}$ over
$J\sim|\log a_0|$ steps and diverge as $a_0\to 0$.

The proposed resolution is \emph{extended extraction}. Instead of extracting
only dimension $\leq 4$ operators, one extracts all gauge-invariant local
operators up to canonical dimension $\dmax = 16$ (for $b=2$). The irrelevant
sector then begins at dimension $\dmax+1 = 17$ and contracts by
\[
  \Lirr \;=\; b^{\,4-(\dmax+1)} \;=\; b^{\,3-\dmax} \;=\; 2^{-13} \;\approx\; 1.2\times 10^{-4}.
\]
With the same rigorous $C_1 \leq 2224$, the effective contraction factor becomes
\[
  \boxed{\;\Lambda \;=\; C_1\,\Lirr \;\leq\; 2224 \cdot 2^{-13} \;\approx\; 0.27 \;<\; 1.\;}
\]
The manuscript presents this as genuine contraction ``with substantial slack.''
The price is tracking a finite-dimensional vector of higher couplings
$\vec g_j = (g_j^{(4)},g_j^{(6)},\dots,g_j^{(\dmax)})$ rather than the single
gauge coupling; the manuscript argues these are perturbatively small and do not
affect the universal one-loop beta-function coefficient
$b_0 = 11/(3\cdot 16\pi^2)$, which is determined entirely by local ultraviolet
information.

\subsection{The recombination constant and the extraction-projection norm}

The constant $C_1$ is itself a product of factors (manuscript Computation~O.7),
\[
  C_1 \;\leq\; C_{\mathrm{fluct}} \cdot C_{\mathrm{cumulant}} \cdot C_{\mathrm{recomb}} \cdot C_{\mathrm{extract}}
  \;\leq\; 1.65 \times 3b^4 \times 14 \times 2 \;\approx\; 139\,b^4,
\]
which for $b=2$ gives $C_1 \leq 139 \times 16 \approx 2224$. The last factor
$C_{\mathrm{extract}} = \|\Pleq\|_{\mathrm{op}}$ is the operator norm of the
extraction projection onto the finite-dimensional space of operators of
dimension $\leq \dmax$. The manuscript bounds it (its Computation~O.7(d)) by
\[
  \|\Pleq\|_{\mathrm{op}} \;\leq\; \sum_{d=0}^{\dmax} \binom{\#\mathrm{indices}}{d}\,(r_1/r_0)^{-d}
  \;\leq\; 1 + O\!\big((r_0/r_1)^{\dmax}\big),
\]
and asserts $\|\Pleq\|_{\mathrm{op}} \leq 2$ for $r_1/r_0 = 1/2$ and $\dmax=16$.
We return to this expression in \S\ref{sec:honest}; it is the load-bearing
input on which the smallness of $C_1$, and hence the contraction, rests.

\section{What we formalized}

The formalization lives in the Lean namespace
\texttt{Mettapedia.QuantumTheory.YangMills}. We summarize the four checked
components: the contraction predicate, the knife-edge arithmetic, the finite
strong-coupling gaps with their finite OS reconstruction, and the
continuum-open gating. All theorem and definition names below are reproduced
verbatim from the corresponding modules.

\subsection{The extended-extraction contraction predicate}

The module \texttt{Mettapedia.QuantumTheory.YangMills.RGBootstrap} models the
route's contraction obligation as a predicate over a recombination constant
$C$, block factor $b$, and extraction depth $\dmax$. The irrelevant scaling
factor is $b^{3-\dmax}$ (an integer power, since the useful cases have
$\dmax>3$), and a contraction is a nonnegative constant whose product with the
irrelevant scale is below $1$.

\begin{lstlisting}
/-- Irrelevant scaling factor `b^(3-dmax)` used by the extended-extraction
discussion.  The exponent is an integer because the useful cases have
`dmax > 3`. -/
def irrelevantScale (b dmax : ℕ) : ℝ :=
  (b : ℝ) ^ (3 - (dmax : ℤ))

/-- The route-level contraction audit for a combinatorial constant and an
irrelevant scaling factor. -/
def HasRGContraction (C scale : ℝ) : Prop :=
  0 ≤ C ∧ 0 ≤ scale ∧ C * scale < 1

/-- The extended-extraction audit specialized to `λ = b^(3-dmax)`. -/
def HasExtendedExtractionContraction (C : ℝ) (b dmax : ℕ) : Prop :=
  HasRGContraction C (irrelevantScale b dmax)
\end{lstlisting}

A threshold lemma in the same module records the exact content of the predicate:
for $b>0$ and $\dmax\geq 3$, the contraction \texttt{HasExtendedExtractionContraction
C b dmax} holds if and only if $0 \leq C < b^{\dmax-3}$. Thus the entire
arithmetic content of the route is the comparison of $C$ with the threshold
$b^{\dmax-3}$.

\subsection{The knife-edge arithmetic}

The module machine-checks that the advertised constant contracts at $\dmax=16$
and that the exact one-step gain is $2224\cdot 2^{-13} = 139/512 < 1/3$:

\begin{lstlisting}
/-- The advertised numerical slack: `2224 * 2^-13 < 1`. -/
theorem rgContraction_2224_two_sixteen :
    HasExtendedExtractionContraction 2224 2 16 := by
  refine ⟨by norm_num, by rw [irrelevantScale_two_sixteen]; norm_num, ?_⟩
  rw [irrelevantScale_two_sixteen]
  norm_num

/-- Exact advertised one-step gain for the extended-extraction constants. -/
theorem rgGain_2224_two_sixteen_eq :
    (2224 : ℝ) * irrelevantScale 2 16 = (139 : ℝ) / 512 := by
  rw [irrelevantScale_two_sixteen]
  norm_num
\end{lstlisting}

Crucially, the same constant does \emph{not} contract at the nearest even
extraction depth below the advertised one. At $\dmax=14$ the threshold is
$2^{11}=2048<2224$, so the contraction is arithmetically false; the module
proves that any contraction at $\dmax=14$ would force $C<2048$, and hence that
$2224$ fails:

\begin{lstlisting}
/-- Any contraction proof at `dmax = 14` must improve the linear recombination
constant below `2048`.  The manuscript's advertised bound `2224` is therefore
insufficient for this lower extraction depth. -/
theorem rgContraction_two_fourteen_forces_constant_lt_2048
    {C : ℝ} (h : HasExtendedExtractionContraction C 2 14) :
    C < 2048 := by
  have hprod : C * irrelevantScale 2 14 < 1 := h.2.2
  rw [irrelevantScale_two_fourteen] at hprod
  nlinarith

/-- Same-constant blocker at the nearest even extraction depth below `16`. -/
theorem not_rgContraction_2224_two_fourteen :
    ¬ HasExtendedExtractionContraction 2224 2 14 := by
  intro h
  have hC : (2224 : ℝ) < 2048 :=
    rgContraction_two_fourteen_forces_constant_lt_2048 h
  norm_num at hC
\end{lstlisting}

This is the knife edge: the choice $\dmax=16$ is exactly what makes the
advertised constant contract, and the manuscript's own table shows
$\dmax\in\{10,12,14\}$ all fail against $C_1\leq 2224$ while $\dmax=16$ succeeds.
The module \texttt{Mettapedia.QuantumTheory.YangMills.RGCrux} packages this into
a route-level statement: a mass-gap route that keeps the advertised constant
$2224$ but uses an even extraction depth strictly below $16$ cannot exist, since
the required contraction is arithmetically false before any OS / clustering
layer is reached.

\begin{lstlisting}
/-- Same-constant lower-extraction certificate in the manuscript's even
canonical-dimension shape: the cutoff is strictly below `16`, is even, and uses
the same advertised recombination bound `2224`. -/
def SameConstantEvenBelowSixteenRGCertificate (dmax : ℕ) : Prop :=
  dmax < 16 ∧ Even dmax ∧ HasExtendedExtractionContraction 2224 2 dmax

/-- No same-constant RG certificate exists for an even manuscript-shape cutoff
strictly below `16`. -/
theorem not_sameConstantEvenBelowSixteenRGCertificate
    {dmax : ℕ} :
    ¬ SameConstantEvenBelowSixteenRGCertificate dmax := by
  intro hcert
  have hdmax : dmax ≤ 14 :=
    even_lt_sixteen_le_fourteen hcert.1 hcert.2.1
  exact (not_rgContraction_2224_two_of_dmax_le_fourteen hdmax) hcert.2.2
\end{lstlisting}

A strengthened companion theorem,
\texttt{not\_atLeast2048EvenBelowSixteenRGCertificate}, shows that any even
lower-extraction repair below $\dmax=16$ would need a genuinely sharper constant
($C<2048$), not merely the advertised one. We stress that this is an honest
delimiter of the arithmetic surface, not a refutation of the route: at the raw
natural number $\dmax=15$ the same constant \emph{does} contract
($2^{12}=4096>2224$), and the module records this explicitly so that the
obstruction is read precisely as an even canonical-dimension statement matching
the manuscript's dimensions $4,6,\dots,\dmax$.

\subsection{Finite strong-coupling spectral gaps and finite OS reconstruction}

The module \texttt{Mettapedia.QuantumTheory.YangMills.FiniteOSReconstruction}
formalizes the finite algebraic core of OS reconstruction: a reflection-positive
finite kernel induces a positive pairing on observables, and a transfer operator
that is self-adjoint for that pairing descends to the pairing quotient and
remains self-adjoint there. On this base, the modules
\texttt{Z2StrongCouplingGap} and \texttt{Z3StrongCouplingGap} build explicit
finite single-link abelian lattice gauge models. For $\mathbb{Z}_2$, the
one-link heat-kernel transfer at heat time $1$ has eigenvalues $1$ and
$\exp(-1)$, so the dimensionless gap is $-\log(\exp(-1)/1)=1$:

\begin{lstlisting}
theorem z2StrongCouplingGap_eq_one :
    z2StrongCouplingGap = 1 := by
  simp [z2StrongCouplingGap, z2StrongCouplingLambdaTwo,
    z2StrongCouplingLambdaOne, z2StrongCouplingRatio, Real.log_exp]

/-- End-to-end finite lattice theorem: the explicit `Z₂` strong-coupling model is
nontrivial, reflection positive, has an OS transfer operator, and has
dimensionless transfer gap at least `1`. -/
theorem z2StrongCoupling_latticeYangMills_massGap_landmark :
    ∃ _ : Z2StrongCouplingGapCanary,
      (∃ x y : Z2OneLinkConfig, x ≠ y) ∧
        OSReflectionPositive z2StrongCouplingKernel ∧
        z2StrongCouplingGap = 1 ∧
        (1 : ℝ) ≤ z2StrongCouplingGap ∧
        0 < z2StrongCouplingGap := by
  exact
    ⟨z2StrongCouplingPositiveCanary,
      z2OneLinkConfig_nontrivial,
      z2StrongCouplingKernel_reflectionPositive,
      z2StrongCouplingGap_eq_one,
      z2StrongCouplingGap_lower_bound_one,
      z2StrongCouplingGap_positive⟩
\end{lstlisting}

The $\mathbb{Z}_3$ module proves the analogous landmark with two mean-zero
transfer eigenvectors and the same unit gap. These theorems are unconditional
and non-vacuous: they exhibit genuine reflection positivity, a genuine finite OS
transfer reconstruction, and a strictly positive spectral gap. They are also
kept rigorously separate from any continuum claim. Each module carries an
explicit scaling diagnostic showing that a finite gap need not survive a naive
continuum schedule. For the heat-time schedule $t(a)=a^2$, the \emph{physical}
gap (dimensionless gap divided by lattice spacing) closes as $a\to0$:

\begin{lstlisting}
/-- A concrete gap-closing schedule for the finite model under the heat-time
scaling `t(a)=a²`: the physical gap can be made smaller than any positive
threshold by taking the lattice spacing small enough. -/
theorem z2QuadraticHeatTime_physicalGap_closes :
    ∀ ε : ℝ, 0 < ε →
      ∃ a : ℝ,
        0 < a ∧ a < ε ∧
          z2HeatTimePhysicalGap a (a ^ 2) < ε := by
  intro ε hε
  refine ⟨ε / 2, by linarith, by linarith, ?_⟩
  have hne : ε / 2 ≠ 0 := by linarith
  rw [z2QuadraticHeatTime_physicalGap_eq hne]
  linarith
\end{lstlisting}

\subsection{The continuum endpoint is gated open}

The module \texttt{Mettapedia.QuantumTheory.YangMills.ProofState} records a
status ledger over the route's nodes. The finite RG / extraction surfaces are
marked \texttt{checked} or \texttt{refuted} as appropriate, but the continuum
mass-gap endpoint carries an explicit \texttt{openGoal} status, gated behind a
constructive-QFT promotion barrier (continuum measure, Euclidean covariance,
reflection positivity, OS reconstruction, and Hamiltonian mass-gap transfer ---
none of which is represented for the RG / extraction route, which supplies only
a finite ledger):

\begin{lstlisting}
/-- The final mass-gap endpoint remains an open goal in this lane. -/
def yangMillsMassGapEndpointNode : YangMillsProofNode where
  id := "yang-mills.mass-gap-endpoint"
  status := .openGoal
  truthValue := ⟨0, 95⟩
  evidence := "currentYangMillsMassGapEndpoint_blockedByConstructiveGate blocks promotion from the audited RG/extraction surfaces to a continuum mass-gap endpoint ..."
  nextObligation := "Connect the RG and extraction surfaces to a genuine continuum mass-gap statement only after the constructive-QFT gate, OS/Hamiltonian transfer dependency packet, challenge-world interface packet, and mass-gap promotion gate clear."

theorem yangMillsMassGapEndpointNode_open :
    yangMillsMassGapEndpointNode.status = .openGoal := by
  rfl
\end{lstlisting}

A companion theorem \texttt{currentYangMillsMassGapEndpoint\_blockedByConstructiveGate}
records that the endpoint is open, the constructive checks are all unrepresented,
and the only supplied interface is the finite RG / extraction ledger. In short,
the formalization does not overclaim: nothing in it asserts a continuum mass gap.

\section{Honest status and scope}\label{sec:honest}

This lane is, by construction, conservative; we now state precisely what is and
is not established.

\paragraph{The contraction constant is an input, not a conclusion.}
Every arithmetic theorem above takes the recombination constant as a
\emph{given real number}. The predicate \texttt{HasExtendedExtractionContraction
C b dmax} is a statement about $C$, $b$, $\dmax$ alone; the checked facts
\texttt{rgContraction\_2224\_two\_sixteen} and
\texttt{not\_rgContraction\_2224\_two\_fourteen} verify the comparison
$C \lessgtr b^{\dmax-3}$ for the specific value $C=2224$. Nothing in the
formalization derives that $2224$ (or any $O(1)\cdot b^4$ bound) is the correct
value of the manuscript's recombination constant. The Lean ledger is therefore
faithful to the route's bookkeeping but agnostic about the route's hardest
analytic claim.

\paragraph{The extraction-projection norm justification is the key
unformalized risk.}
The smallness of $C_1$ rests on the factor
$C_{\mathrm{extract}} = \|\Pleq\|_{\mathrm{op}}$, which the manuscript bounds by
\[
  \|\Pleq\|_{\mathrm{op}}
  \;\leq\; \sum_{d=0}^{\dmax} \binom{\#\mathrm{indices}}{d}\,(r_1/r_0)^{-d}
  \;\leq\; 1 + O\!\big((r_0/r_1)^{\dmax}\big),
\]
asserting the value $\leq 2$ for $r_1/r_0 = 1/2$, $\dmax = 16$. We record
neutrally a tension internal to this written expression: with $r_1/r_0 = 1/2$
the summand $(r_1/r_0)^{-d} = 2^{d}$ \emph{grows} geometrically in $d$, and the
binomial weights only enlarge it further, so the written partial sum is an
\emph{increasing} function of $\dmax$ that grows at least like $2^{\dmax}$; the
asserted closed form $1 + O\big((r_0/r_1)^{\dmax}\big) = 1 + O(2^{-\dmax})$
instead behaves as if the summand were $(r_1/r_0)^{d} = 2^{-d}$, i.e. as if it
\emph{decayed}. We do not adjudicate which reading is intended. We observe only
that this single step decides whether $C_1$ stays $O(1)\cdot b^4 \approx 2224$ or
instead grows with the extraction depth $\dmax$ --- and that, because the
irrelevant factor is $\Lirr = 2^{3-\dmax}$ while $C_{\mathrm{extract}}$ enters
$C_1$ multiplicatively, a $C_{\mathrm{extract}}$ that grew like $2^{\dmax}$ would
overwhelm the very $2^{-13}$ contraction the scheme relies on. This is precisely
the part not captured in Lean. We flag it as the key unformalized risk; we do
not assert that the constant is wrong, only that its smallness is not
self-evident from the manuscript's own stated expression and is not yet
machine-checked.

\paragraph{The finite gaps are real but Clay-irrelevant.}
The $\mathbb{Z}_2$ and $\mathbb{Z}_3$ spectral gaps are honest theorems, but
they are abelian, single-link, strong-coupling, and finite. They exemplify the
finite OS-reconstruction mechanism; they are not lattice approximations to a
continuum $\SUN$ theory, and the accompanying scaling diagnostics show that the
physical gap can close under a naive continuum schedule. They establish the
machinery, not the target.

\paragraph{The continuum mass gap is correctly open.}
The continuum endpoint is gated \texttt{openGoal} with no claimed support, and
the constructive-QFT prerequisites are all marked unrepresented. The
formalization makes no claim, in either direction, about the existence of a
continuum Yang--Mills mass gap.

\section{What remains}

To turn this lane into evidence about the route's central claim rather than its
bookkeeping, the decisive next step is to formalize, or independently bound, the
extraction-projection operator norm $\|\Pleq\|_{\mathrm{op}}$ as a function of
$\dmax$, and thereby determine whether the recombination constant $C_1$ is
genuinely $O(1)\cdot b^4$ or grows with the extraction depth. Until that input is
established, the checked contraction arithmetic --- correct as it is --- rests on
an assumed value. Beyond it lie the genuinely infinite-dimensional layers the
route requires and the present ledger leaves unrepresented: the KP polymer
expansion and tree-decay clustering at the stopping scale, the glue /
back-propagation transfer of clustering from the coarse to the fine scale, the
thermodynamic and continuum OS limits, rotation restoration, the
transfer-matrix spectral gap, and the nontriviality witness. Each is an
analytic theorem, not an arithmetic comparison, and none is currently
machine-checked. The Yang--Mills mass gap remains open.

\paragraph{Lean modules.}
The results reported here are in
\texttt{Mettapedia.QuantumTheory.YangMills.RGBootstrap} (the contraction
predicate and knife-edge arithmetic),
\texttt{Mettapedia.QuantumTheory.YangMills.RGCrux} (the even lower-extraction
route obstruction),
\texttt{Mettapedia.QuantumTheory.YangMills.FiniteOSReconstruction},
\texttt{Mettapedia.QuantumTheory.YangMills.Z2StrongCouplingGap}, and
\texttt{Mettapedia.QuantumTheory.YangMills.Z3StrongCouplingGap} (the finite OS
reconstruction and strong-coupling gaps), and
\texttt{Mettapedia.QuantumTheory.YangMills.ProofState} (the continuum-open
gating). The primary mathematical source is Ben Goertzel's manuscript
\emph{Yang--Mills Mass Gap in 4D: RG Bootstrap, Polymer Expansion, OS
Reconstruction} (v14).

\end{document}
